\chapter{最短的学习回路}

\section{不要等到下周才知道桥塌了}

这本书的材料很杂：课堂讲授、学生汇报、教师培训、游戏原型、网页工具。把它们串起来的不是 AI，也不是桥，而是一个反复出现的动作：\textbf{做一次，看见结果，改一下，再解释。}

2026 年 5 月 22 日，我用最朴素的话说过它：

\begin{shiyu-inline}
``我做出来一个操作，马上获得这个东西行还是不行，好还是不好。然后我跟这个回应，再去做调试。''

\hfill ——游戏化教学分享，2026 年 5 月 22 日 00:13:32
\end{shiyu-inline}

这句话后来被原稿不断拔高，几乎变成了游戏解决一切。其实它只解决第一件事：缩短动作和结果之间的距离。

\section{四步，不是两步}

完整回路有四步：

\begin{enumerate}
  \item \textbf{尝试：}学生先作出预测、模型或初稿；
  \item \textbf{反馈：}系统、同伴、教师或真实约束指出差距；
  \item \textbf{修改：}学生据此改变作品，而不是只读评语；
  \item \textbf{解释：}学生说明改了什么、为什么、还不确定什么。
\end{enumerate}

只做“尝试—反馈”，学生可能看完分数就走。只做“生成—修改”，修改可能完全由 AI 完成。最后的解释让教师看见判断是否真的回到了学生手里。

反馈研究把重点从“教师给了什么”转向“学习者拿它做了什么”。Carless 把跨任务持续使用反馈称为反馈螺旋\cite{carless2019spirals}。这与课程项目里的版本迭代非常接近：上一版的问题，应该进入下一版，而不是随成绩一起封存。

\section{谁来给反馈}

桥本身可以反馈：车开过去，结构变形或破坏。测试用例可以反馈：代码在边界输入下报错。同伴可以反馈：你的解释哪里听不懂。AI 可以反馈：版本之间变了什么、哪些评论重复。教师则把时间留给模型假设、专业取舍和学生自己还没有察觉的盲点。

反馈来源越多，不代表越好。最关键的是每一种反馈只做自己擅长的事，并且学生有机会行动。

\section{教师真正困难的工作}

设计回路比讲解更麻烦。你要决定在哪一步让学生卡一下，什么时候给提示，什么错误可以安全发生，什么错误必须立即制止。你还要忍住，不在学生刚开始想的时候就把答案说出来。

更现实的问题是：一个班几十份作品、几轮迭代，教师会不会被拖成一台永远改不完的机器？下一章，我们不谈“教师不可替代”这种宽慰，直接谈哪些活该交给 AI，哪些活教师必须留给自己。
